\section{模型检验、误差分析与评价}
\subsection{分层检验思路}
\paragraph{从公共物理量到全局方案}
验证应沿模型依赖链展开。先确认航段经过像元、作业高度和能耗单位，再核对单架次计算，之后检验全局资源与通信。若公共地形函数存在偏差，即使输出表之间完全一致，也不能据此证明任务满足50 m净空。

单架次检查包含四项：货箱质量与体积、随交付递减的剩余载荷、逐站交接时刻、累计能耗与返航SOC。全方案再检查每箱只出现一次、全部硬截止满足、实体及能源区间不重叠。问题三在此基础上对每段轨迹给出连续通信依据；问题四则重建实际保障关系，核对分组前后的固定任务和时刻。

\paragraph{检验指标的判定层次}
本文把检验结果分为“物理量正确、单任务可行、全局约束可行、目标值可回算”四层。第一层检查距离、高度、时间、能量和损耗的单位及边界；第二层检查每个运输或中继任务的容量、SOC和链路；第三层检查80箱覆盖、硬时限和资源并发；第四层从逐项结果重新汇总架次数、能耗、完工时间和资源峰值。只有四层同时通过，汇总结果才被写入正文。

这种层次能够区分两类错误：若单架次能耗不正确，应回到公共物理模型；若单架次均可行但资源区间重叠，则问题出在调度；若明细正确而摘要数字不同，则属于汇总或表述不一致。定位层次后再修正，可以避免用调整最终数字掩盖底层计算问题。

\paragraph{可行性、守恒关系与交叉复核}
\paragraph{四问关键检验}
表~\ref{tab:model-check}汇总四个问题的关键检验。覆盖守恒按货箱编号计数，要求80箱各出现一次；资源检验采用左闭右开区间扫描，结束事件先于同一时刻的开始事件处理；能量检验直接由逐航段能耗回算返航SOC，而不是使用汇总值反推。问题三同时使用连续区间证书与固定步长网格，两种方法回答“保守证明”和“独立抽查”两个不同问题。

\begin{table}[htbp]
\centering
\caption{四个问题的模型检验结果}\label{tab:model-check}
\small
\begin{tabularx}{\textwidth}{>{\centering\arraybackslash}p{0.8cm}>{\raggedright\arraybackslash}p{2.8cm}YY}
\toprule
问题 & 检验对象 & 判定准则 & 结果\\
\midrule
Q1 & 组批覆盖与返航能量 & 80箱各出现一次；每架次返航SOC不低于下限 & 通过；18架次，最小返航SOC为22.789\%\\
Q2 & 交付、时限与双资源占用 & 唯一覆盖；31个硬时限箱零违约；机体与电池区间无重叠 & 全部通过\\
Q3 & 连续通信、链路与中继资源 & 未解析区间为0；0.5~s网格中断点为0；中继机与组件无冲突 & 22条轨迹均通过，65532个网格点中断为0\\
Q4 & 分区完备性与资源公式 & 枚举非空带标签候选并核对无标签等价类；缺口、冗余逐类型复算 & $K=2$共1023个；$K=3$共57002个带标签候选，对应28501个无标签等价类\\
\bottomrule
\end{tabularx}
\end{table}

\paragraph{基线、等价公式与反向回算}
问题二还保留26架次硬箱专送基线作为对照。22架次混装方案不仅保持硬违约为0，而且$WTD$、完工时间、能耗和架次数同时下降，因此改进不是由放松硬约束获得。问题四的最小资源数由区间峰值和贪心复用两种等价口径核对；在固定时序与同型资源可互换的假设下，两者得到相同配置数。

能量方面，对每个运输架次同时保存逐航段能耗之和和返航SOC，二者应满足$SOC^{end}=1-E/E^{use}$；对问题三，运输能耗与中继能耗分别汇总后再相加，得到79.097126~kWh。时间方面，完工时间直接取最后返航事件，而不是充电结束事件；资源恢复时间则另外从全部充电事件取最大值。两种时间口径分别回算，可防止把“任务完成”和“系统完全恢复”混为一谈。

问题四还进行任务归属反向检查：对每个分组，把组内服务区映射回运输架次，再从实际通信分段映射到中继架次；所有原运输任务必须至少归属一个组，同一原子块不能跨组。对共享中继任务按独立执行规则在相关组中复制完整任务，再计算资源峰值。由此保证资源缺口确实来自分区独立性，而非漏记或错误切分任务。

\subsection{敏感性实验与解释边界}
\paragraph{实验设计与重算范围}
表\ref{tab:validation-plan}区分本次已经重算的敏感性实验和仍作为扩展方向的项目。返航余量实验重新计算安全载荷与问题一组批，地面延误实验在既定问题二事件链上逐级传播，链路保守余量实验则从连续区间证书的最小裕量出发判断原方案是否仍被认证。能耗系数扰动和中继任务局部裁剪会改变候选方案本身，本文不把它们写成已经完成的稳健性结论。
\begin{table}[htbp]\centering
\caption{敏感性实验与扩展检验范围}\label{tab:validation-plan}
\begin{tabularx}{\linewidth}{>{\raggedright\arraybackslash}p{2.3cm}YY}\toprule
实验对象 & 参数变化与重算范围 & 本文结论或状态\\\midrule
返航余量 & 10\%、15\%、20\%、25\%、30\%；重算安全载荷与组批 & 已完成：架次数依次为18、18、18、19、20\\
地面操作延误 & 0--1200~s，每60~s传播至问题二时序 & 已完成：硬逾期始终为0，最小硬时限余量降至1055.086~s\\
链路保守余量 & 额外收紧0、0.1、0.2、0.5、1、2~dB & 已完成：0.1~dB时22架次仍通过，0.2~dB时仅16架次通过旧证书\\
能耗不确定性 & 水平能耗系数及爬升效率分别扰动，先独立再联合比较 & 扩展项目：需重新判断能量可行性并重排方案\\
分区执行方式 & 主方案保持完整中继任务；局部时段裁剪作为另行假设 & 扩展项目：需重算分类型配置、缺口、冗余与均衡\\
\bottomrule\end{tabularx}\end{table}

返航余量从10\%提高到20\%时，最优组批仍保持18架次、59.240~kWh和32798.500~s，最小返航SOC为22.789\%，说明20\%基准在当前离散装箱结构下尚未触发重组。余量增至25\%后，安全载荷集合收缩使架次数升至19，能耗升至61.190~kWh；增至30\%时进一步升至20架次和67.358~kWh。该阶梯变化说明，安全载荷是连续变化的，但完整货箱的可行子集会在临界点发生离散跳变，因此不能仅报告某个机型的连续载荷曲线。

\paragraph{地面延误传播}
地面延误不能用“全部交付时刻统一加常数”替代，因为一次延误会沿同一实体机和共享电池的后续事件传播。本次将0--1200~s附加延误按60~s步长作用于既定调度的相关事件链；所有情景下硬逾期箱数均为0，最小硬时限余量由2255.086~s线性降至1055.086~s。因测试上界仍未越过硬时限，结论只能表述为方案对1200~s以内的该类延误保持硬约束可行，不能外推为任意延误下均可行。一般物资的软时限代价和资源完全恢复时间仍会随事件推迟而增加，实际执行时应在滚动调度中更新。

\paragraph{链路保守余量}
通信鲁棒性应区分“旧保守证据不再足以保证”与“实际链路中断”。基准连续证书的最小下界裕量为0.1397~dB；额外收紧0.1~dB后，22个运输架次仍由原证书覆盖，最差调整后裕量为0.0397~dB；收紧0.2~dB后只有16个架次仍被旧证书认证，收紧0.5、1和2~dB时分别为8、4和3个。这里的“未认证”不等于已经证明断链，而表示原点位和原区间证书不足以保证更严格条件。若要构造相应鲁棒方案，必须重新执行区间验证、选点和调度，并比较新增中继能耗与联合完工时间。

三组敏感性输出与主结果一同纳入统一复核记录。它们分别回答离散组批对安全余量的响应、固定调度对事件延误的硬时限余量，以及既定通信证书对附加链路保守量的承受范围。三者没有把尚未重算的能耗系数扰动或局部中继裁剪混入已完成结论，从而保持“改变参数后重算什么”与“当前证据能证明什么”的边界。

从决策角度看，返航余量实验给出装备安全要求提高时的运输代价阶梯，延误实验给出既定资源编排能够吸收的时间缓冲，链路实验则揭示当前通信方案的最小证书余量很窄。三类结果应分别用于组批预案、现场调度和中继点位加固，不能相互替代。

\subsection{误差来源、数值精度与影响}
\paragraph{误差分类与控制方式}
模型输出的误差不能只理解为浮点舍入。对本题影响更大的，是DEM空间分辨率、传播模型简化和启发式搜索边界。表~\ref{tab:error-analysis}按“来源—当前控制—剩余影响”给出误差审计。

\begin{table}[htbp]
\centering
\caption{主要误差来源及其控制方式}\label{tab:error-analysis}
\small
\begin{tabularx}{\textwidth}{>{\raggedright\arraybackslash}p{2.5cm}YY}
\toprule
误差来源 & 当前控制与可量化精度 & 对结论的影响与处置\\
\midrule
安全载荷求根 & 二分迭代60次，质量区间远小于输出的$10^{-3}$~kg精度 & 数值误差可忽略；临界箱组仍按未四舍五入值判定\\
DEM离散 & 沿航段遍历30~m DEM经过像元并取最高地形，统一增加50~m净空 & 对给定DEM是保守的；真实小尺度障碍和高程测量误差未被附件描述，现场需以更高精度地形复核\\
时间与事件计算 & 内部以秒和双精度计算，约束比较容差为$10^{-9}$量级；表中时间通常保留$10^{-3}$~s & 展示舍入误差不超过0.0005~s，不参与可行性判断\\
连续通信认证 & 未能直接证明的区间递归二分至0.1~s并按不可用处理；另以0.5~s网格独立抽查 & 数值检查无中断，但最小保守余量仅0.1397~dB，传播模型误差比浮点误差更重要\\
能耗参数 & 采用附件确定值，水平能耗与爬升能耗逐段累计 & 未包含风场、电池老化和温度效应；返航余量敏感性显示25\%起会触发重新组批\\
组合搜索 & Q1集合划分与Q4固定方案分区为穷举/动态规划；Q2、Q3为确定性启发式 & Q2、Q3只证明当前方案可行且优于所列基线，不证明全局最优\\
\bottomrule
\end{tabularx}
\end{table}

\paragraph{空间误差的传播路径}
DEM误差首先改变航段最高地形，继而同时影响巡航海拔、爬升时间、爬升能耗和视线遮挡。若某一高程被低估，问题一可能高估安全载荷，问题二可能低估交付时间和能耗，问题三还可能把遮挡链路误判为可用，最终传播到问题四的资源需求。因此，地形不是只影响一张地图的外部数据，而是贯穿四问的公共误差源。

本文对给定30~m DEM采用经过像元最大值加50~m的保守口径，但这种保守性只针对栅格表示本身。分辨率以下的树木、建筑和临时障碍不在附件中，无法由数值算法消除。现场部署应使用更高分辨率地形或实测航线复核，并重新运行能耗和通信模型。

\paragraph{时间误差与事件传播}
准备、装载、交接、建链和周转时间以题目给定值进入事件链。单个事件偏差不仅影响当前任务，还会通过实体机或能源资源的可用时刻影响后续任务。若多个任务连续复用同一电池，早期延误会累积传播；若换用另一块已满电电池，则传播路径可能中断。因此，误差不能简单地对所有交付时刻加同一个常数。

本文地面延误实验沿实际资源前后继关系传播，能够反映这种结构，但尚未覆盖多类事件同时随机波动。更完整的随机稳健分析可为准备、交接和建链时间设定分布，采用蒙特卡洛仿真统计硬时限失效概率；在缺少分布依据时，本文不虚构置信区间。

\paragraph{算法误差与最优性边界}
问题一在每个服务区完整枚举箱组并动态规划，问题四对固定问题三方案完整枚举合法分区，两者的算法误差主要来自输入和数值计算。问题二只比较两套确定性构造，问题三使用离线筛选后固化的点位与波次；它们证明当前方案可行且优于所列基线，不代表全局最优。相应结果必须用“可行方案、相对基线改善、当前可行上界”等措辞，而不能用未证明的“全局最短”或“绝对最优”。

由此可见，计算舍入不是主要风险。问题一在20\%返航余量下最小实际返航SOC为22.789\%，数值容差不会改变可行性；问题三最小保守链路余量只有0.1397~dB，对未建模附加损耗更敏感。因此通信结论必须限定为题设传播模型和给定DEM下成立，现场部署应预留至少0.1~dB以外的工程裕量，或重新选点并认证。

\subsection{模型评价与适用边界}
\paragraph{模型优点}
模型的主要优点有三点。第一，运输、能源、通信和分区共用同一事件时间轴，避免把机体返航等同于电池恢复；第二，完整货箱、逐站剩余载荷和实际通信分段均保留到输出层，结果能够追溯到具体对象；第三，连续区间证书、敏感性重算和基线对照共同约束结论边界，而非只报告一个目标值。载荷相关无人机路径与空中通信的既有研究也强调能量、路径和链路条件的耦合\textsuperscript{\cite{Dorling2017,Zeng2016}}。

\paragraph{模型局限}
局限同样明确。问题二只在硬箱专送基线和同目的地混装方案之间择优，问题三又在冻结运输路线后执行固化的中继点位与波次，因此22个运输架次和11个中继架次均不具备全局最优性证明；联合完成时间6.572~h应理解为当前可行上界。30~m DEM与二值遮挡附加损耗无法表达树木、建筑、降雨和小尺度多径，0.1397~dB的最小链路余量提示通信方案偏紧。问题四的结论只适用于问题三的固定任务和绝对时序；若允许各组重新排程，资源缺口与均衡性需重新计算。

\paragraph{适用场景}
因此，本模型适合标准场景下的方案生成、资源核算和预案比较，不应直接替代现场飞行许可、无线电勘测与气象评估。实际执行时宜采用滚动重调度：更新天气、道路与链路观测后，重新计算剩余任务的能量、时限和中继点位。

\subsection{模型特点与可推广方向}
\paragraph{统一事件时间轴}
传统做法容易把路线优化、机体排程、电池充电和中继保障分开计算，最后才发现时间上无法拼接。本文把准备开始、起飞、逐箱交付、返航、充电、建链和周转全部表示为事件，并通过资源前后继关系传播。这一表示使问题二的运输调度可以直接扩展到问题三，也使问题四能够在不重排时刻的条件下准确计算独立资源峰值。

\paragraph{连续认证与离散复核分工}
通信方案先用离散采样快速筛选候选，再用连续区间证书判定可行，最后以更密的0.5~s网格作独立复核。三者分别承担效率、证明和查错功能，避免把“采样未发现中断”等同于“连续时间必然可用”。这一思路也可推广到地形净空、动态禁飞区或其他连续轨迹约束。

\paragraph{从可执行性到滚动优化}
当前模型输入为确定场景，输出完整任务与资源时序。实际灾害现场可采用滚动时域方式：每完成一批任务，就用最新天气、SOC、链路观测和需求变化更新剩余任务，只冻结已经开始的区间，再重新运行组批、资源排程和连续认证。这样既保留当前模型的可解释约束，又能适应现场状态变化。

\paragraph{图表与数值的一致性检查}
\paragraph{图、表与明细数据的对应}
正文结果由同一份逐架次、逐箱和通信分段数据生成，汇总指标均可逐项回算。问题一说明容量瓶颈，问题二说明交付与资源时序，问题三说明连续保障与代价，问题四说明独立分区造成的资源变化。相同指标的比较图使用一致量纲和色标；示意图不承担数值验证功能。

附录\ref{app:outputs}列出结果表保留的关键字段。这些字段同时服务计算复核与图表解释，使“图上看到的现象”能够定位到具体架次、货箱或时间区间。

本研究在信息检索、赛题理解与任务拆解、建模讨论、辅助编程、数据质量检查、图表和论文排版等环节使用了生成式人工智能工具；工具信息、代表性输入及输出处理方式见附录\ref{app:ai-usage}。
